%======================================================================
\section{Robustness}\label{sec:robust}
%======================================================================

This section asks whether the base-case split of the AT1 spread into about $180$bp priced
in Tier~2 and $50$bp of AT1-specific premium survives three objections a practitioner would
raise: that the call rule is unrealistic, that the attribution depends on the order of
features, and that a single flat resolution hazard is too crude for a perpetual. A fourth
objection, that the split depends on the weakly identified speed of capital mean reversion,
is treated in Section~\ref{sec:uncert}. It is the one the split does not fully survive: the
Tier~2-priced part is robust to it, the AT1-specific premium is not. All computations are on
the valuation date, on the same random numbers as the base case.

\subsection{Reputational calls}

The base call rule \eqref{eq:call} with $\delta=0$ calls only when refinancing is cheaper.
European national champions have almost always called their AT1s at the first opportunity,
and market participants quote call probabilities well above those of
Table~\ref{tab:callext}. We introduce a reputational tolerance: the issuer also calls when its
reset margin is up to $\delta$ \emph{below} the new-issue spread. For each $\delta$ we
recalibrate the hazard of each bond.

\begin{table}[H]
\centering
\begin{tabular}{lcccc}
\toprule
Tolerance $\delta$ (bp) & 0 & 25 & 50 & 100 \\
\midrule
$\Pr(\text{called at }T_1)$: 6.875\% NC33 & 49\% & 57\% & 65\% & 77\% \\
\phantom{$\Pr(\text{called at }T_1)$:} 7.20\% NC36 & 52\% & 60\% & 67\% & 79\% \\
\phantom{$\Pr(\text{called at }T_1)$:} 7.45\% NC35 & 58\% & 66\% & 72\% & 81\% \\
Implied hazard, 6.875\% (\%/yr) & 2.21 & 2.22 & 2.23 & 2.25 \\
Implied hazard, 7.20\% (\%/yr) & 2.18 & 2.19 & 2.20 & 2.21 \\
Implied hazard, 7.45\% (\%/yr) & 2.19 & 2.19 & 2.19 & 2.19 \\
Extension component, 6.875\% (bp) & 15 & 14 & 13 & 11 \\
Extension component, 7.20\% (bp) & 9 & 9 & 8 & 7 \\
Extension component, 7.45\% (bp) & 6 & 6 & 6 & 6 \\
Tier~2 share of the AT1 hazard & 78\% & 78\% & 78\% & 78\% \\
\bottomrule
\end{tabular}
\caption{Reputational call tolerance. The extension component is the yield-to-call spread of
the model price minus that of the same bond with a certain call at $T_1$.}
\label{tab:reputation}
\end{table}

A tolerance of $100$bp raises the probability of a call at the first reset to $77$--$81\%$,
the range practitioners quote, and lowers the extension component of the low-margin bond
from $15$ to $11$bp. The implied hazard moves by at most $4$bp, and the Tier~2 share of the
hazard does not move. The reason is that the market price pins down the total value of the
bond: a more generous call rule raises the value of the call leg slightly, and the hazard
adjusts to absorb it. The extension risk of the low-margin bond, and its negative convexity,
are therefore partly a feature of the economic call rule, while the split between Tier~2-priced
and AT1-specific risk is not.

\subsection{Order of attribution}

Table~\ref{tab:order} reports four attributions of the spread for each bond: the base order
(R, S, M, C, E), the nested Shapley value of Section~\ref{sec:attribution}, the full Shapley
value over all $120$ orders, and the reverse order (E, C, M, S, R).

\begin{table}[H]
\centering\small
\begin{tabular}{llcccc}
\toprule
Bond & Component & Base order & Nested Shapley & Full Shapley & Reverse order \\
\midrule
\multirow{5}{*}{6.875\% NC33} & Priced in Tier~2 & 180.7 & 180.8 & 221.1 & 259.1 \\
 & AT1-specific & 51.6 & 51.6 & 62.8 & 77.0 \\
 & MDA & 27.3 & 36.1 & 35.5 & 43.3 \\
 & Conversion & 4.1 & 5.5 & 6.4 & 11.7 \\
 & Extension & 14.7 & 4.5 & $-47.3$ & $-112.6$ \\
\midrule
\multirow{5}{*}{7.20\% NC36} & Priced in Tier~2 & 180.8 & 180.9 & 206.9 & 231.4 \\
 & AT1-specific & 48.8 & 48.8 & 55.7 & 64.4 \\
 & MDA & 30.6 & 36.0 & 35.1 & 39.3 \\
 & Conversion & 4.6 & 6.0 & 6.8 & 11.0 \\
 & Extension & 9.1 & 2.4 & $-30.4$ & $-72.0$ \\
\midrule
\multirow{5}{*}{7.45\% NC35} & Priced in Tier~2 & 180.0 & 180.0 & 203.6 & 225.9 \\
 & AT1-specific & 48.9 & 48.8 & 55.1 & 62.9 \\
 & MDA & 30.4 & 35.8 & 35.0 & 39.2 \\
 & Conversion & 4.0 & 5.3 & 6.1 & 10.1 \\
 & Extension & 6.0 & $-0.6$ & $-30.4$ & $-68.8$ \\
\bottomrule
\end{tabular}
\caption{Attribution of the yield-to-call spread (bp) under four methods. Totals are
$278.5$, $274.0$ and $269.3$bp in every column.}
\label{tab:order}
\end{table}

Three conclusions follow.

\emph{First}, the split of the hazard block is order-free. In every attribution that prices
the hazard before the contract features, the base order and the nested Shapley value, the
Tier~2-priced and AT1-specific components agree to within $0.2$bp. This is the near-linearity
of Proposition~\ref{prop:hazardsplit}: on the relevant range the yield-to-call spread is
almost linear in the hazard level, so it does not matter in which order $\lres$ and $\lspec$
are switched on.

\emph{Second}, the attribution among the contract features is order-dependent. Their total
is the same in the base order and in the nested value ($40$--$46$bp depending on the bond),
but its split is not: symmetric treatment raises the MDA component to about $36$bp and
reduces extension to a few basis points. The two features interact through the capital
condition of the call rule. The states in which the call is blocked because CET1 is below the
threshold are the states in which the MDA cuts coupons; once the MDA is priced, the
cash flows of the bond in those states are already reduced, so the incremental cost of not
being called is smaller, and conversely. The honest statement is that the contract features
cost $40$--$46$bp, of which MDA is $27$--$36$bp, conversion $4$--$6$bp and extension between
$-1$ and $15$bp.

\emph{Third}, orders that switch extension on before the hazard produce economically
meaningless results. Without any hazard, the bond is risk-free, and not calling a risk-free
bond that pays its reset margin over Treasuries is a gift to the investor: extension takes a
large negative value ($-69$ to $-113$bp in the reverse order) and the hazard components absorb
the difference. This is the inconsistency the state-dependent hazard is designed to avoid
(Section~\ref{sec:hazard}); it shows why the full Shapley value, which averages over such orders,
is the wrong benchmark for AT1 attribution.

\subsection{Term structure of the resolution hazard}

A single flat $\lres$, calibrated to a bond with a 2034 call, is applied to AT1 cash flows
that run for decades if the bond is extended. With the two Tier~2 bonds we bootstrap a
two-piece resolution hazard as in Section~\ref{sec:piecewise}.

\begin{table}[H]
\centering
\begin{tabular}{lccc}
\toprule
 & 6.875\% NC33 & 7.20\% NC36 & 7.45\% NC35 \\
\midrule
$\lambda_1$ to August 2030 (\%/yr) & 1.40 & 1.40 & 1.40 \\
$\lambda_2$ after August 2030 (\%/yr) & 2.10 & 2.10 & 2.10 \\
Flat $\lres$ (\%/yr) & 1.72 & 1.72 & 1.72 \\
$\lspec$, flat $\lres$ (\%/yr) & 0.49 & 0.46 & 0.47 \\
$\lspec$, two-piece $\lres$ (\%/yr) & 0.42 & 0.36 & 0.39 \\
Tier~2-priced component, flat (bp) & 180.7 & 180.8 & 180.0 \\
Tier~2-priced component, two-piece (bp) & 176.4 & 184.6 & 181.2 \\
\bottomrule
\end{tabular}
\caption{Two-piece resolution hazard bootstrapped from the Tier~2 2035NC30 (call August 2030)
and the Tier~2 2035NC34 (call November 2034), recovery $20\%$.}
\label{tab:termstructure}
\end{table}

The resolution hazard is upward sloping: $1.40\%$ a year to 2030 and $2.10\%$ after. With
the steeper curve, more of each AT1's hazard is attributed to what Tier~2 prices, and the
AT1-specific hazard falls from $0.47\%$ to $0.39\%$ on average. The Tier~2-priced component
of the spread moves by less than $5$bp, in different directions for different bonds
depending on the timing of their cash flows. A term structure therefore strengthens, rather
than weakens, the conclusion that most of the AT1 spread is priced in Tier~2. A caveat
applies to the short Tier~2: it trades well below par with a reset in 2030, and its price
may contain a convexity or tax effect that a bullet-to-call treatment ignores.